% ID: SEL-D0055-OB-E04
% ESTADO: cotejada. Segunda pasada visual y borrador compilado revisados.
% Fuente: Selectividad/sel_2011_matematicas_aplicadas/122A-EXAMEN 6.pdf
% SHA-256: aebcbfddd4d5db5f5a10b4662e9e12b2ee362bbf574b751bccd6d78594766ad3
% Página: 2; opción B; ejercicio original 4.

El peso neto de las tabletas de chocolate de una determinada marca es una variable aleatoria Normal con media $\mu$ y desviación típica 7 gramos. Se sabe que 36 tabletas, elegidas al azar, han dado un peso total de 5274 gramos.
\begin{enumerate}[label=\alph*)]
\item \textbf{(1.25 puntos)} Calcule un intervalo con un nivel de confianza del 94\% para la media $\mu$.
\item \textbf{(1.25 puntos)} Con el mismo nivel de confianza, \textquestiondown{}cuántas tabletas, como mínimo, habrá que tomar como muestra para que la amplitud del intervalo que se obtenga sea, como máximo, de 3 gramos?
\end{enumerate}
